\section{The first cause and the main theorem}\label{sec:firstcause}

Every reduction of \cref{sec:reduction} presses inward, toward a single open node --- the last
Step-00 obligation --- narrowed moreover to the scale $A \ge 5$. There is nowhere further inward to
go. The present section turns around and looks outward, to answer three questions. Why is it
legitimate to accept the last node as an external first cause rather than as a hidden theorem? What
does such an architecture prove \emph{unconditionally}, with no axiom at all? And, deepest,
\emph{why} can the node not be closed from inside --- why is knowledge of its cause unattainable for
an observer living inside the system? The answer to the third question is our main theorem.

\subsection{The cause cannot be derived from inside}

With an open node one may proceed in two ways: leave it open (honest, but the story stops), or
accept it (and everything downstream closes). The second path is legitimate under exactly one
condition --- the acceptance must be \emph{visible}: one may declare an explicit assumption and
trace what rests on it, but one must never pass the unproven off as proven. It turns out there is no
freedom in \emph{how} to accept. To accept the node from inside --- to exhibit an internal
self-grounding of it --- is impossible, and this is a theorem.

An internal proof that crosses its own boundary in order to ground that boundary would, step by step
under machine check, build a forbidden perpetual engine; since no engine exists (the acyclicity of
the descent rank, \cref{sec:engine}), no such self-grounding can exist either.

\begin{theorem}\label{thm:no_internalSelfDerivation_step00CausalClosure}
\green{}\ The causal boundary does not certify itself: there is no internal derivation that grounds
the boundary by crossing it. \leantag{no\_internalSelfDerivation\_step00CausalClosure}
\end{theorem}

The same statement in the language of a beginning: a self-caused first frame --- an internalised
event $0 \to 1$ --- cannot exist in a stable, engine-free architecture.

\begin{theorem}\label{thm:no_internalisedHorizonBoundary}
\green{}\ An internalised causal horizon --- a first frame that grounds itself from within --- does
not exist. \leantag{no\_internalisedHorizonBoundary}
\end{theorem}

\begin{remark}
By itself the prohibition of self-certification is close to tautological: proving $P$ while crossing
the boundary of $P$ collapses into $P \wedge \lnot P$. What is substantive is the
\emph{identification of form} --- that the attempt at self-grounding is already a burnt engine
construction, not some new impossibility. Physically it is the perpetual engine of the first kind:
self-grounding would extract the work of derivation from nothing. The descent rank is a conserved
quantity with a strict floor, and the loop breaks against that floor. Hence the foundation must be
brought in from outside.
\end{remark}

\subsection{An axiom accepted intentionally: the event \texorpdfstring{$0\to1$}{0->1}}

Since the cause cannot be laid from inside, we lay it from outside, deliberately, as the single axiom
of the entire development, the \emph{first cause} \axiom{}. Its type is the structure of the event
$0 \to 1$, the transition from non-being to the first causal frame. Two of its five fields are pure
markers carrying only $\mathrm{True}$ --- the singularity $0$ (where there is no internal language
yet) and the first frame $1$. The remaining three are the substantive causal boundaries, one per
branch of the programme.

\begin{definition}\label{def:step00FirstCause}
The first cause carries exactly three causal boundaries: the \emph{twin boundary} (\code{causalBoundary}),
which is precisely the open node of \cref{sec:reduction}; the \emph{Riemann boundary}
(\code{riemannBoundary}), the law of manifestation of a zero (\cref{sec:riemann}); and the
\emph{Navier--Stokes boundary} (\code{nsBoundary}), the gate law of the energy balance
(\cref{sec:ns}). All mathematical weight lives in these three fields; the two markers assert nothing.
\leantag{Step00FirstCause}
\end{definition}

The framing changes the provenance of the result, not its strength: the decree is exactly the sum of
what is put into it.

\begin{theorem}\label{thm:step00FirstCause_iff_causalClosure}
\green{}\ The first cause is equivalent to the conjunction of its three boundaries: the twin node,
the Riemann law, and the Navier--Stokes gate. \leantag{step00FirstCause\_iff\_causalClosure}
\end{theorem}

The axiom is locked in a \emph{quarantine module}, and the quarantine is a machine procedure: a
separate verifier marks every declaration depending on \axiom{} as axiom-tainted, and none has leaked
into the green line (\cref{sec:status}). The former closure axiom is now a theorem obtained by
projecting the first cause onto its twin boundary. Along the twin branch it yields the infinitude of
the lower twin centres --- \emph{conditionally}.

\begin{theorem}\label{thm:twinLowersInfinite_from_step00CausalClosure}
\yellow{}\ Under the accepted twin boundary, the lower twin centres are infinite.
\leantag{twinLowersInfinite\_from\_step00CausalClosure}
\end{theorem}

\emph{This is a conditional derivation, not a proof of the twin-prime conjecture.} The decree is
moreover no weaker than its consequence: it exhibits a twin centre above every threshold. Accepting
the axiom \emph{is} accepting the twins; it creates no free strength. The axiom-free remainder is
honest to the point of tautology.

\begin{theorem}\label{thm:nonAxiomaticRemainingObligation_iff_lastStep00Obligation}
\green{}\ The axiom-free remaining obligation is exactly the old open node \red{}: no wrapper lowers
the open content below the node, and the twin-prime conjecture remains a \code{sorry} that does not
close through the quarantine. \leantag{nonAxiomaticRemainingObligation\_iff\_lastStep00Obligation}
\end{theorem}

\subsection{Epistemics: the cause exists, cannot be known, and knowledge breaks everything}

Three things can be said of the internal knowledge of the cause. It \emph{exists} --- that is the
axiom \axiom{}, \yellow{}, accepted intentionally. It \emph{cannot be known} --- and this is a
theorem, \green{}, not a caveat.

\begin{theorem}\label{thm:cause_unknowable}
\green{}\ Internal knowledge of the first cause is impossible. \leantag{cause\_unknowable}
\end{theorem}

\emph{Proof idea.} To know the cause from inside would be to derive it from inside, which builds a
perpetual engine; no engine exists. The unknowability of the cause is not a weakness but a
conservation law.

Knowledge, were it possible, would be explosive: through the impossible engine it entails both that
the twins are finite and that they are infinite --- both \emph{ex falso}. The substantive,
non-explosive residue is the dichotomy \emph{either the cause is unknowable, or the twins are
finite}, whose left disjunct is \cref{thm:cause_unknowable}: we are always in the left world. The
load-bearing content is not the explosion but an axiom-free lemma --- the heart of the causal line
--- in which the hypothesis ``no engines'' is consumed substantively rather than vacuously.

\begin{theorem}\label{thm:twins_infinite_of_noEngine_and_boundary}
\green{}\ If no perpetual engine exists and the causal boundary is accepted, then the twin primes are
infinite. \leantag{twins\_infinite\_of\_noEngine\_and\_boundary}
\end{theorem}

\begin{proof}
Suppose the twins are finite; then some last twin centre $M_0$ has none above it. From this finite
boundary one builds an infinite family of generated flows; the finite semantic key is forced by
pigeonhole to collide; and from the collision a concrete witness engine is assembled. That witness is
exactly what the hypothesis ``no perpetual engine'' forbids --- not an empty explosion but a direct
hit. The hypothesis is consumed substantively.
\end{proof}

Instantiating this lemma by the descent rank's absence of engines and by the boundary from the axiom
gives the twins by decree.

\begin{theorem}\label{thm:twins_because_unknowable}
\yellow{}\ The twin primes are infinite, by a derivation that visibly passes through the same ``no
engines'' fact that yields unknowability. \leantag{twins\_because\_unknowable}
\end{theorem}

\begin{remark}
``The twins are infinite because the cause cannot be known'' means \emph{through a common cause and a
load-bearing lemma}, not ``unknowability alone proves infinitude''. All the weight of \emph{existence}
still rests on the accepted boundary: unknowability without the boundary yields no twins, or they
would already be proven.
\end{remark}

\subsection{The two walls}

There is a second, independent reason the twins escape the observer --- cognitive rather than causal,
and unconditional (its core relies only on propositional extensionality). A finite system of level
$A$ can recognise a twin only through the finite data it sees, and that data is shared by an entire
sieve-equivalence class of centres. Hence it recognises a twin $B$ only if its \emph{whole} class
consists of twins; a twin sitting in a class that also contains a non-twin is invisible in principle;
and if the classes remain mixed arbitrarily far out, infinitude cannot be certified from inside at
all. That the classes of a \emph{non-trivial} view actually mix at concrete levels is an arithmetic
input \red{}; the structural theorems and the myopic (trivial-view) instantiation are unconditional.
The two walls close into one statement.

\begin{theorem}\label{thm:two_walls_one_nature}
\green{}\ From inside a finite view a twin whose sieve-equivalence class also contains a non-twin
cannot be recognised; and from inside the system the first cause cannot be known. The two
impossibilities hold conjointly. \leantag{two\_walls\_one\_nature}
\end{theorem}

Both walls are one thing: a finite system cannot cross its own boundary. The finite view of level $A$
is a cognitive horizon, just as the strict order of traversal is a causal one.

\subsection{The main theorem: higher-energy incompatibility}

Everything gathers into a single statement, entirely in the green core, that glues the causal wall to
the cognitive wall.

\begin{theorem}\label{thm:higherEnergyIncompatibility_main}
\green{}\ \emph{(no axioms)} Five faces of one incompatibility hold conjointly:
\begin{enumerate}[label=\textup{(\arabic*)},leftmargin=2.2em,itemsep=1pt]
\item knowledge of the first cause from inside builds a perpetual engine;
\item therefore the first cause is unknowable from inside;
\item a finite view recognises a twin only if its entire sieve-equivalence class consists of twins;
\item under persistent mixing of the classes, infinitude cannot be certified from inside;
\item the load-bearing face: this very incompatibility --- ``no engines $=$ the cause cannot be
  known'' --- together with the accepted boundary, entails that the twin primes are infinite.
\end{enumerate}
\leantag{higherEnergyIncompatibility\_main}
\end{theorem}

Faces (1)--(4) are the two walls side by side; face (5) is the bridge from them to the conclusion.
Read energetically: knowing from inside costs a perpetual engine, which does not exist, while the
infinitude of the twins is external knowledge, paid for by the first cause. Instantiating face (5) by
decree gives the corollary --- yellow, above a green core.

\begin{corollary}\label{thm:higherEnergyIncompatibility_twins}
\yellow{}\ The twin primes are infinite. \emph{This is not a proof of the twin-prime conjecture}; the
conjecture remains a \code{sorry} that does not close through the quarantine.
\leantag{higherEnergyIncompatibility\_twins}
\end{corollary}

\begin{remark}
The name is not decoration. Landauer's principle makes information physical --- erasing a bit costs
at least $kT\ln 2$, and Maxwell's demon respects the second law precisely because its knowledge must
be paid for. Our theorem is a discrete sibling: knowing the cause from inside would cost infinite work
(a perpetual engine), hence is forbidden by the same conservation law; the infinitude of the twins is
not free, its price being knowledge carried in from outside, that is, an axiom. This draws out three
worlds. If the node is true, the axiom is superfluous and the twins are in fact proven; if the node is
independent, the theory is consistent but this cannot be known from inside; if the node is refutable,
the quarantine is inconsistent and a tripwire fires at once, devaluing every tainted declaration. Both
internal resolutions cost a perpetual engine, which is why the boundary can be accepted only from
outside.
\end{remark}
